\chapter{Giới thiệu đề tài}
\label{chap:gioi-thieu}

\section{Bối cảnh và lý do chọn đề tài}
Tại quầy đồ uống, khách hàng thường cần xem menu, hỏi giá theo kích cỡ,
chọn topping rồi chờ nhân viên ghi nhận và thanh toán. Khi nhiều đơn được
đặt cùng lúc, thao tác thủ công dễ làm sai tùy chọn, nhầm tổng tiền hoặc
nhận đơn vượt lượng nguyên liệu còn có thể phục vụ.

BrewLite là ứng dụng web hỗ trợ khách tự chọn đồ uống và thanh toán
không tiền mặt theo luồng đặt món tại quầy. Một đơn hàng giữ được các tùy
chọn và giá tại thời điểm đặt; hệ thống kiểm tra nguyên liệu trước khi
xác nhận và cung cấp thông tin để khách theo dõi đơn. Đây cũng là bài toán
để nhóm thực hành transaction, xử lý đồng thời và chống ghi nhận thanh
toán trùng trong phạm vi môn Công nghệ phần mềm.

\section{Mục tiêu}
\begin{itemize}
    \item Cung cấp luồng khách hàng từ xem menu, tùy chỉnh đồ uống,
    quản lý giỏ, đăng nhập, tạo đơn đến thanh toán và xem lại đơn.
    \item Xác nhận giá ở backend, lưu bản chụp thông tin món và bảo vệ
    dữ liệu đơn hàng theo tài khoản đăng nhập.
    \item Quản lý khả năng phục vụ theo nguyên liệu và công thức;
    giữ, sử dụng hoặc hoàn nguyên liệu theo vòng đời đơn hàng.
    \item Mô phỏng thanh toán Ví/Thẻ, hỗ trợ thử lại và xử lý một yêu cầu
    thanh toán lặp mà không tạo tác động tài chính trùng.
    \item Tổ chức phát triển theo Scrum, có phân công, tiêu chí chấp nhận,
    mã nguồn, kiểm thử và tài liệu bàn giao để đối chiếu.
\end{itemize}

\section{Phạm vi và đối tượng sử dụng}
\begin{table}[H]
\centering\small
\caption{Đối tượng và nhu cầu sử dụng BrewLite}
\label{tab:actors}
\begin{tabularx}{\textwidth}{@{}L{0.23\textwidth}Y@{}}
\toprule
\textbf{Đối tượng} & \textbf{Nhu cầu và ranh giới} \\
\midrule
Khách chưa đăng nhập & Xem menu, chi tiết và chuẩn bị giỏ hàng; cần đăng nhập trước khi gửi đơn. \\
Khách hàng (CUSTOMER) & Đặt món, thanh toán, xem xác nhận, lịch sử và chi tiết đơn của mình. \\
Nhân viên (BARISTA) & Nhu cầu nhận đơn, pha chế và giao món; tài khoản/role đã có, giao diện và API vận hành quầy còn trong kế hoạch. \\
Mock Payment Gateway & Tác nhân tích hợp trả kết quả thanh toán giả lập; không thực hiện giao dịch tiền thật. \\
Nhóm phát triển, giảng viên & Theo dõi backlog, đánh giá chức năng, chất lượng và kết quả thực hiện. \\
\bottomrule
\end{tabularx}
\end{table}

Phạm vi sản phẩm tập trung vào đồ uống tại một cửa hàng. Các nội dung như
vận chuyển, đặt bàn, quản lý chuỗi cửa hàng, thanh toán qua ngân hàng thật
và hoàn tiền thật không thuộc phần hiện thực đang báo cáo. Coupon và
loyalty là nghiệp vụ hỗ trợ luồng mua hàng; sự tồn tại của bảng dữ liệu
không đồng nghĩa đã có đầy đủ giao diện quản trị.

\section{Công nghệ sử dụng}
\begin{table}[H]
\centering\small
\caption{Công nghệ đối chiếu từ cấu hình và mã nguồn}
\label{tab:stack}
\begin{tabularx}{\textwidth}{@{}L{0.21\textwidth}Y@{}}
\toprule
\textbf{Thành phần} & \textbf{Công nghệ và vai trò} \\
\midrule
Frontend & Next.js, React, TypeScript, Tailwind CSS; các trang trong \sourcefile{apps/web/src/app/}. \\
Trạng thái giao diện & TanStack Query cho dữ liệu API; giỏ hàng dùng store tự xây dựng với \code{useSyncExternalStore} và \code{localStorage}. \\
Backend & NestJS, TypeScript, REST API; tách module Auth, Catalog, Orders, Inventory, Payments và Promotions. \\
Lưu trữ & PostgreSQL, Prisma ORM, migration có phiên bản; tiền lưu bằng số nguyên VND. \\
Xác thực & JWT, Passport, bcryptjs; kiểm tra role và quyền sở hữu tài nguyên. \\
Kiểm thử & Vitest ở backend, bộ kiểm thử Node ở frontend và Playwright cho trình duyệt. \\
Hạ tầng & Docker Compose, GitHub Actions; Node.js 24 theo cấu hình của hai ứng dụng. \\
\bottomrule
\end{tabularx}
\end{table}

\section{Cơ sở xây dựng báo cáo}
Báo cáo tham khảo cách trình bày User Story và Acceptance Criteria trong
\sourcefile{assigment1_tham_khao.pdf}, đặc biệt sáu câu chuyện của luồng
khách hàng~\cite{reference-assignment}. Nội dung công nghệ, chức năng và
phân công được đối chiếu lại với mã nguồn, schema, Product Backlog và hồ
sơ nhóm~\cite{source-code,backlog,assignment}. Các mô tả cũ như dùng
Zustand hoặc lưu tồn trực tiếp theo thành phẩm không được dùng để mô tả
hiện trạng.

Báo cáo gồm sáu chương: giới thiệu; phân tích yêu cầu; thiết kế hệ thống;
hiện thực; kiểm thử; tổ chức thực hiện và kết luận. Mã US, FR và NFR lần
lượt nhận diện User Story, yêu cầu chức năng và yêu cầu phi chức năng,
phục vụ truy vết giữa nhu cầu, cách hiện thực và kiểm chứng.
